\usepackage{ctex}
\usepackage{fontspec}
\setCJKmainfont{Noto Serif SC}
\setCJKsansfont{Noto Sans SC}
\setmonofont{DejaVu Sans Mono}[Scale=0.90]

\usepackage[a4paper,top=2.35cm,bottom=2.35cm,left=2.25cm,right=2.25cm]{geometry}
\usepackage{xcolor}
\usepackage{hyperref}
\usepackage{tikz}

\definecolor{dynasty}{HTML}{7A3326}
\definecolor{ink}{HTML}{332B25}
\definecolor{timeline}{HTML}{9A6B3E}
\definecolor{paper}{HTML}{FFF9EE}
\definecolor{source}{HTML}{665A4B}
\hypersetup{
  colorlinks=true,
  linkcolor=dynasty,
  citecolor=source,
  urlcolor=timeline,
  bookmarksnumbered=true,
  plainpages=false,
  hypertexnames=false,
  pdftitle={中国大事编年·第二卷：汉、三国、两晋与南北朝}
}
\color{ink}

% Keep event links compatible with the other chronicle volumes.
\newcommand{\sourcebook}[1]{《#1》}
\newcommand{\eventid}[1]{\textcolor{dynasty}{\footnotesize\textsf{〔#1〕}}}
\newcommand{\eventanchor}[2]{\hypertarget{evt:#1}{\eventid{#2}}}
\newcommand{\eventref}[2]{\hyperlink{evt:#1}{\textcolor{dynasty}{〔#2〕}}}

% Use these only after the target event anchor or figure label exists.
\newcommand{\sourceindexentry}[2]{\item[\sourcebook{#1}] #2}
\newcommand{\eventindexentry}[2]{\item \eventref{#1}{#2}}
\newcommand{\figureindexentry}[2]{\item \hyperref[#1]{“#2”}}
